% Generated draft: requires mathematical and source review.
\documentclass[11pt]{article}
\usepackage[margin=25mm]{geometry}
\usepackage{fontspec}
\setmainfont{Latin Modern Roman}
\usepackage{amsmath,amssymb,mathtools,unicode-math}
\setmathfont{Latin Modern Math}
\usepackage{enumitem,needspace,xurl,hyperref}
\setlength{\parindent}{0pt}
\setlength{\parskip}{5pt}
\setlength{\emergencystretch}{3em}
\newcommand{\N}{\mathbb{N}}
\newcommand{\Z}{\mathbb{Z}}
\newcommand{\Q}{\mathbb{Q}}
\newcommand{\R}{\mathbb{R}}
\newcommand{\C}{\mathbb{C}}
\newcommand{\F}{\mathbb{F}}
\newcommand{\A}{\mathbb{A}}
\newcommand{\PP}{\mathbb{P}}
\newcommand{\E}{\mathbb{E}}
\newcommand{\eps}{\varepsilon}
\newcommand{\dd}{\,\mathrm{d}}
\begin{document}
\section*{UnsolvedMath ID: 30000439}
\textbf{Source problem code:} catalogue entry
\subsection{Definitions and mathematical statement}
\textbf{Embedding-gap question.} Can the minimum dimensions for piecewise-linear embeddability and linear embeddability of a simplicial complex differ by more than one? For a nonempty finite simplicial complex \(K\), define\[e_{\mathrm{lin}}(K)=\min\{m:|K|\text{ embeds in }\R^m\text{ by a map affine on each original simplex}\},\] and\[e_{\mathrm{PL}}(K)=\min\{m:|K|\text{ has a piecewise-linear embedding in }\R^m\}.\] In the second definition a subdivision may be used to make the map affine on its simplices. Thus \(e_{\mathrm{PL}}(K)\le e_{\mathrm{lin}}(K)\), and the requested existence statement is\[\exists K:\quad e_{\mathrm{lin}}(K)-e_{\mathrm{PL}}(K)\ge2.\]\par A linear embedding in this simplicial sense assigns points to the original vertices, sends simplices to their affine convex hulls, and is globally injective, so distinct simplex images intersect only in the image of their common face. It is not required to be a single linear transformation of an externally fixed realization. The combinatorial triangulation \(K\) must remain fixed when evaluating \(e_{\mathrm{lin}}\): allowing arbitrary subdivisions in both definitions would erase the intended distinction. Conversely, replacing PL embeddings by arbitrary topological embeddings would change the first minimum. The source does not explicitly say finite, locally finite, or which infinite-complex category is allowed; the displayed finite-complex formulation is the standard interpretation rather than an added universal claim about infinite complexes.\par The supplied triage checked 2026-08-21 labels the problem open. It reports Brehm--Sarkaria examples with a one-dimensional gap and says its search found no result establishing or excluding a gap greater than one. The actual examples and their dimensions are not reproduced. The question asks only whether a gap at least two occurs, not whether gaps can be arbitrarily large. This dated assessment requires independent verification.
\subsection{Short English statement}
Does some fixed triangulated complex need at least two more ambient dimensions for an embedding affine on its original simplices than for a PL embedding after subdivision? Keeping the original triangulation fixed in the linear minimum is essential. The supplied triage reports examples with gap one.
\subsection{Sources}
\begin{itemize}\item Local UnsolvedMath registry, permanent ID 30000439, code catalogue entry; \emph{Gaps Between Piecewise-Linear and Linear Embedding Dimensions}, Oberwolfach Reports Open Problems.\item \emph{Discrete Differential Geometry (2006)}, Oberwolfach Reports, DOI 10.4171/owr/2006/12, pp. 653--728. \url{https://doi.org/10.4171/owr/2006/12}. Proposed year 2006; no proposer named.\item Supplied literature label \emph{original OWR problem}: \url{https://doi.org/10.4171/owr/2006/12}.\item Supplied literature label \emph{Brehm--Sarkaria}: \url{https://archive.mpim-bonn.mpg.de/id/eprint/1946/1/preprint_1992_52.pdf}. This is a supplied reference only; it was not downloaded.\item Literature triage checked 2026-08-21, status \texttt{open}; record updated 2026-09-16. Statement status \texttt{exact}, with source checking reported in the verification field, not independently checked here.\end{itemize}
\end{document}
